\section{Architettura di rete}

\subsection{Protocolli e comunicazione}
Tutte le richieste fra mobile app e backend avvengono via \textbf{HTTP/1.1}
(in produzione su \textbf{HTTPS} con TLS terminato da un reverse proxy).
Il formato di scambio è \textbf{JSON} con codifica UTF-8.

I client autenticati allegano il token di accesso nell'header
\code{Authorization: Bearer <jwt>}; le pubblicazioni anonime usano invece
l'header custom \code{X-Publish-Token: <prefix>:<nonce>} (vedi
\cref{sec:sicurezza}).

\subsection{Metodi HTTP e codici di stato}
\begin{table}[H]
    \centering
    \begin{tabularx}{\textwidth}{@{}llX@{}}
        \toprule
        \textbf{Metodo} & \textbf{Uso} & \textbf{Esempio} \\
        \midrule
        GET & lettura risorse, idempotente, cacheabile & \endpoint{GET}{/api/get/3} \\
        POST & creazione di risorse o richieste di azione & \endpoint{POST}{/api/publish} \\
        PUT & aggiornamento completo di una risorsa esistente & \endpoint{PUT}{/api/songs/3} \\
        DELETE & rimozione di una risorsa & \endpoint{DELETE}{/api/songs/3} \\
        \bottomrule
    \end{tabularx}
    \caption{Verbi HTTP utilizzati.}
\end{table}

\begin{table}[H]
    \centering
    \begin{tabularx}{\textwidth}{@{}llX@{}}
        \toprule
        \textbf{Codice} & \textbf{Significato} & \textbf{Quando viene restituito} \\
        \midrule
        200 & OK & Risorsa restituita o operazione riuscita \\
        201 & Created & Nuova risorsa creata (es. publish, register) \\
        400 & Bad Request & Body non valido o parametri obbligatori mancanti \\
        401 & Unauthorized & Credenziali assenti, malformate o non valide \\
        403 & Forbidden & PoW fallita o operazione non autorizzata per l'utente \\
        404 & Not Found & Risorsa inesistente \\
        409 & Conflict & Username/email già in uso, oppure token PoW già consumato \\
        413 & Payload Too Large & Body oltre \code{MAX\_CONTENT\_LENGTH} \\
        422 & Unprocessable Entity & Validazione semantica fallita (formati, lunghezze) \\
        500 & Internal Server Error & Errore non gestito lato server \\
        \bottomrule
    \end{tabularx}
    \caption{Codici HTTP restituiti dal sistema.}
\end{table}

\subsection{Schema di errore}
Tutti gli errori delle API \code{/api/*} seguono lo schema LRCLIB:
\begin{lstlisting}[language=json,caption={Esempio di risposta di errore.}]
{
    "statusCode": 422,
    "error": "PublishError",
    "message": "Il campo 'plainLyrics' eccede la lunghezza massima"
}
\end{lstlisting}

\subsection{Misure di sicurezza di rete}
\begin{itemize}
    \item \textbf{TLS / HTTPS:} in produzione il traffico è cifrato end-to-end
    tramite reverse proxy (\emph{nginx} o \emph{Caddy}) con certificato
    \emph{Let's Encrypt} e header \code{Strict-Transport-Security} attivo.
    \item \textbf{Firewall:} la macchina di backend espone solo la porta 443
    pubblicamente; la porta 5000 di Flask resta in \code{localhost}.
    \item \textbf{Rate limiting:} previsto su \code{/auth/login},
    \code{/auth/register}, \code{/api/publish} e \code{/api/request-challenge}
    via \emph{Flask-Limiter} (es. 5 req/min per IP), per mitigare bruteforce e
    spam (\emph{deferred to production}).
    \item \textbf{VPN per accessi amministrativi:} accesso al server di
    backend (SSH, console di amministrazione) ristretto a una rete VPN
    aziendale.
    \item \textbf{CORS:} di default aperto in sviluppo; in produzione
    whitelistare le origin del frontend ufficiale.
\end{itemize}
